\documentclass[10pt,a4paper]{article}
\usepackage[margin=2.2cm]{geometry}
\usepackage{fontspec}
\usepackage{kotex}
\setmainfont{Apple SD Gothic Neo}
\setsansfont{Apple SD Gothic Neo}
\usepackage{amsmath,amssymb}
\usepackage{graphicx}
\usepackage{booktabs}
\usepackage{caption}
\usepackage[dvipsnames]{xcolor}
\usepackage[colorlinks=true,linkcolor=MidnightBlue,citecolor=MidnightBlue,
            urlcolor=MidnightBlue]{hyperref}
\usepackage{titlesec}
\titleformat{\section}{\large\bfseries}{\thesection}{0.6em}{}
\titleformat{\subsection}{\normalsize\bfseries}{\thesubsection}{0.5em}{}
\captionsetup{font=small,labelfont=bf}
\setlength{\parskip}{0.35em}
\setlength{\parindent}{0pt}

% 주장 상자 — 모든 주장은 반증 조건을 함께 적는다
\usepackage{mdframed}
\newmdenv[linewidth=0.4pt,linecolor=gray!60,backgroundcolor=gray!5,
          innertopmargin=6pt,innerbottommargin=6pt,skipabove=8pt,skipbelow=8pt]{claimbox}
\newcommand{\claim}[2]{\begin{claimbox}\textbf{주장.} #1\par\smallskip
  \textcolor{gray!70}{\textbf{반증 조건.} #2}\end{claimbox}}

\title{\bfseries 빈 열은 공짜다\\[2pt]
\large 차원 비용은 열 수가 아니라 \emph{갈릴 값}의 수다}
\author{월드모델 실험실 · 노트 663}
\date{2026-08-06}
\begin{document}\maketitle

\section{한 줄}

값이 아예 없는 열을 하나 넣었다. 판 차가 \textbf{정확히 $0.0000$}, 씨앗 SE
$0.0$, 열한 도메인 전부 $0.0$ 이다. 둘 넣어도 같다. \textbf{내 예측이 틀렸고,
그 덕에 차원 비용의 정의가 좁아졌다.}

\section{왜 쟀나}

노트 662 가 지역 방문자 축으로 \textbf{신호 몫 $+0.0064$}(문턱 $0.0045$ 초과)를
얻었는데 \textbf{순효과는 $+0.0016$} 이었다. 위약이 $-0.0048$ 이므로 열값이 신호의
3/4 를 먹는다. 그 열값이 어디서 오나.

단서가 있었다 --- \verb|_idol| 은 추가 축을 \textbf{열한 도메인 전부}의
설계행렬에 붙이고, 자료가 없는 도메인에는 \verb|0.5| 상수와 마스크 $0$ 을 준다.
방문자 축은 자료가 둘에만 있으므로 \textbf{아홉 도메인이 아무것도 안 담은 열을
하나씩 받는다}.

\claim{그 아홉 개의 빈 열이 열값의 대부분이다.}{빈 열만 넣었을 때 판이 안
움직이면 이 주장이 틀렸다.}

\section{쟀다 --- 정확히 0.0000}

\begin{center}\small
\begin{tabular}{lrrr}\toprule
열 & 판 차 & 씨앗 SE & 양수 \\ \midrule
빈 열 1개 & $\mathbf{0.0000}$ & $0.0$ & 0/12 \\
빈 열 2개 & $\mathbf{0.0000}$ & $0.0$ & 0/12 \\ \bottomrule
\end{tabular}\end{center}

배선을 먼저 확인했다 --- \textbf{11개 도메인에 붙었고 마스크${>}0$ 합이 $0$}
이다. 열이 \emph{있으면서} 비었다. 노트 652 에서 무효한 대조가 정확히 $0.0000$ 을
냈고 그것이 ``뺀 축이 애초에 없었다''의 서명이었으므로, 이번엔 그 확인을 먼저
했다. \textbf{그러므로 이 $0.0000$ 은 무효가 아니라 참값이다.}

\section{무엇이 확정됐나}

\verb|0.5| 상수 · 마스크 $0$ 열은 나무가 \textbf{쪼갤 값이 없어} 후보 분할에서
아예 빠진다. 그래서 열 수가 늘어도 비용이 $0$ 이다.

\claim{차원 비용은 ``열이 많다''가 아니라 ``갈릴 수 있는 값이 많다''에서 온다.}
{값이 있는 행 수를 고정한 채 열 수만 늘렸을 때 판이 내려가면 이 주장이 틀렸다.}

노트 641 이 ``얇은 도메인에서 열 하나당 $\rho$ $0.02\sim0.025$''라고 적었다.
그것은 \textbf{값을 담은 열}에 대한 것이었고, 이 측정이 그 조항을 좁힌다.

\begin{center}\small
\begin{tabular}{lrl}\toprule
열의 종류 & 판 차 & \\ \midrule
빈 열(어디에도 값 없음) & $0.0000$ & 나무가 무시한다 \\
섞은 열(310행에 값) & $-0.0048$ & 잡음을 학습한다 \\
진짜 열(310행에 값) & $+0.0016$ & 신호가 잡음을 넘는다 \\ \bottomrule
\end{tabular}\end{center}

\section{그래서 고칠 것이 없다 --- 그것도 결과다}

예측이 맞았으면 다음 수는 ``자료 있는 도메인에만 열을 붙이게 하네스를 고친다''
였다. \textbf{고칠 것이 없다.} 하네스는 이미 최적이고, 노트 662 의 ``행을
늘리는 것이 유일한 길''이 \emph{소극적 추론}에서 \textbf{측정된 결론}이 됐다 ---
열값은 값이 있는 행 수에만 달렸고 그 행이 $310$ 뿐이다.

\section{덤 --- 상수 하나를 얻었다}

앞으로 \textbf{부분 덮음 축을 넣을 때 ``덮이지 않은 도메인 때문에''를 비용으로
계산하지 않는다.} 계산할 것은 오직 \emph{값이 있는 행 수}다.

이것이 노트 646 의 관측도 정리해 준다 --- 전국 축이 \textbf{덮음 $0$ 인 아이돌에서
$+0.0567$ 을 벌었다}. 빈 열의 값이 $0$ 이므로 그 이득은 \emph{빈 열 때문이 아니고},
그 축이 다른 도메인의 분산을 흡수한 결과다(노트 413 기제). \textbf{두 기제가
이 측정으로 갈렸다.}

\begin{figure}[h]\centering
\includegraphics[width=\textwidth]{figs/f1.pdf}
\caption{(a) 빈 열을 하나·둘 넣어도 판 차가 정확히 $0$ 이다($\times$ 는 내
예측). (b) 비용은 열이 \emph{갈릴 값}을 담을 때만 생긴다. (c) 그래서 노트 641 의
조항이 좁아진다.}
\end{figure}

\end{document}
